\section{An immersion test for entry and recognition}\label{sec:immersion}

There is another way to specify a metric on the given smooth domain.
Every smooth homotopy four-sphere is stably parallelizable
\cite[Theorem 3.1]{KervaireMilnor1963Groups}. In rank five the
stabilization is already in the stable range over a four-dimensional
complex. Hence \(TM\oplus\mathbb R\) is trivial, and the immersion
theorem gives a smooth \(F:M\to\mathbb R^5\)
\cite[\S6]{Hirsch1959Immersions}. Give \(M\) the induced metric
\(g\), write \(A\) for its shape operator, \(H=\operatorname{tr} A\),
\(A^\circ=A-(H/4)I\), and set
\[
 J=\int_M(-\det A)_+\,dV_g,\qquad B=\|A^\circ\|_{L^4(g)}^2 .
\]
Self-intersections are counted on the smooth domain.

\begin{theorem}[Immersed recognition and entry]\label{thm:immersed-entry}
Each of the following conditions on one immersion implies
\(M\cong S^4\) in the smooth category:
\[
       J<8\pi^2/3,\qquad E_+(g)\le64\pi^2/9,\qquad
                     B\le(16\sqrt2/3)\pi .
\]
The conditions are alternatives. The energy condition holds equivalently for any representative
conformal to the induced metric. It requires no Yamabe sign assumption.
On a hypothetical nonstandard type, an immersion with
\(E_+(g)<8\pi^2\) has exactly two index-two critical points in
more than three quarters of all height directions (normalized
spherical measure).
\end{theorem}

\begin{proof}
Choose the unit normal so that its Gauss map \(n:M\to S^4\)
has oriented pullback tangent bundle \(TM\). Euler classes give
\(\deg n=1\), and its signed Jacobian is \(\det A\).
Put \(\omega=|S^4|=8\pi^2/3\).
For almost every \(u\in S^4\), the height \(\langle F,u\rangle\)
is Morse, with critical points \(n^{-1}(\{u,-u\})\). The area
formula gives the Chern--Lashof average
\cite{ChernLashof1957Curvature}
\[
 \frac1{\omega}\int_{S^4}\#\operatorname{Crit}\langle F,u\rangle\,du
  =\frac2{\omega}\int_M |\det A|\,dV_g=2+4J/\omega .
\]
If \(J<\omega\), some Morse height has at most four critical
points, since each count is even. A two-critical-point height
gives a double of two smooth disks, which is standard because any
boundary diffeomorphism is isotopic to an orthogonal map
\cite{Hatcher1983Smale}. With four critical points the ranks in
the Morse complex give one extra pair in adjacent indices.
An extra \(0,1\) pair joins two distinct disks and cancels; duality
handles \(3,4\). For a middle pair reverse the function if
necessary to have one 2- and one 3-handle. The disk with those
two handles has \(H_2=0\), since only the final 4-handle is missing.
It is a standard disk by the handle consequence of Property R
\cite[Corollary 8.6]{Gabai1987FoliationsIII}. Equivalently the
2-handle has zero framing on a knot whose surgery is \(S^1\times S^2\);
Property R makes the knot trivial. Its trace is \(S^2\times D^2\); gluing it to the dual \(S^1\times D^3\) with any boundary marking gives \(S^4\) \cite[Theorem A]{LaudenbachPoenaru1972Handlebodies}. The final
disk gluing gives the same standard sphere. Hence every immersion
of a nonstandard candidate must have \(J\ge\omega\).

In fact every Morse function on a nonstandard candidate has at least
two index-two critical points. Reorder handles by index. The union
of its 0- and 1-handles is connected: no higher attaching region
can connect different components. Amalgamate the disks and 1-handles
of a spanning tree of this connected graph into one disk. The
remaining 1-handles reattach to its actual boundary. Do the same
for the disjoint dual 0- and 1-handlebody. Transport the intervening
framed 2-handle attachments with these boundary identifications.
Their number \(m_2\) is unchanged, and Euler characteristic now says
\(m_1+m_3=m_2\). If \(m_2\le1\), the preceding disk and single
middle-pair arguments give the standard sphere.

Let \(U\) be the open set where the shape operator has two eigenvalues
of each sign, counted with multiplicity. At a Morse critical point
the height Hessian is the second fundamental form up to sign. Reversing
the normal preserves index two. Thus, writing \(m_2(u)\) for the
original height count, the indexed area formula gives
\begin{equation}\label{eq:middle-mass}
 \frac1\omega\int_{S^4}m_2(u)\,du=\frac2\omega\int_U\det A\,dV_g .
\end{equation}
In particular the integral on the right before the factor \(2/\omega\)
is at least \(\omega\) on a nonstandard candidate. The comparison
of the Weyl tensor on middle signatures with the indexed height integral
is due in general dimension to Onti and Vlachos
\cite[\S\S2--3]{OntiVlachos2017Hypersurfaces}. Here the smooth
four-dimensional handle conclusion gives the lower count even though
the middle Betti number vanishes.

In a principal orthonormal frame let \(k_1,\ldots,k_4\) be the
principal curvatures. The Gauss equation gives the scalar curvature
\(R=3H^2/4-|A^\circ|^2\). On both chiral Weyl blocks its eigenvalues
are \((s_j-\bar s)/2\), where \(s_j=k_1 k_l+k_m k_n\) for the
three pairings and \(\bar s\) is their mean. Their squared Hilbert--Schmidt norm is
\(\frac1{12}\sum_{j<k}(s_j-s_k)^2\). On \(U\), write the curvatures
as \((-a,-b,c,d)\) with \(a,b,c,d>0\). The positive pair sum
\(ab+cd\) is at least \(2\sqrt{\det A}\) and each negative pair
sum is at most \(-2\sqrt{\det A}\). Therefore
\begin{equation}\label{eq:middle-weyl}
             |W^\pm|^2_{\rm HS}\ge\tfrac83\det A\quad\text{on }U .
\end{equation}
Equality requires \(ab=cd,\ ac=bd,\ ad=bc\), so \(a=b=c=d\).
Equality throughout a nonempty open subset of \(U\) is impossible.
Indeed the two smooth principal distributions there have rank two
and eigenvalues \(k,-k\) with \(k>0\). Codazzi, tested on orthogonal
vectors in either distribution, gives the directional derivative
of its eigenvalue zero along that distribution. Together they give
\(dk=0\). For \(X\) in the first distribution and \(Y\) in the
second, \((\nabla_X A)Y=(\nabla_Y A)X\) now lies in both
distributions and vanishes. Metric compatibility then shows each
distribution is parallel in all directions. A mixed sectional
curvature would be zero, contrary to the Gauss value \(-k^2\).
Thus when \(\int_U\det A>0\) the integrated comparison is strict.
Equations~\eqref{eq:middle-mass}--\eqref{eq:middle-weyl} force
\(E_+(g)>(8/3)\omega=64\pi^2/9\) on a nonstandard candidate.
More precisely the excess is the sum of three nonnegative terms,
\[
 E_+(g)-\tfrac83\omega =
 E_+(g;M\setminus U)
 +\int_U\bigl(|W^+|^2_{\rm HS}-\tfrac83\det A\bigr)\,dV_g
 +\tfrac43\omega\left(\frac1\omega\int_{S^4}m_2(u)\,du-2\right).
\]
If \(E_+(g)<8\pi^2\), the mean index-two count is less than \(9/4\).
All regular counts are at least two, so the proportion at least
three is less than \(1/4\). This proves the height assertion.
In particular four mutually orthogonal directions can be chosen
with this count simultaneously: the first four vectors of a random
orthonormal frame each have uniform marginal direction, and the
union of the four failure events has probability less than one.
After the two spanning-tree reductions such a two-middle-handle
description has \((m_1,m_3)=(0,2),(1,1)\) or \((2,0)\).

If the four trace-free principal curvatures are
\[
 (x+y+z,\ x-y-z,\ -x+y-z,\ -x-y+z)/2,
\]
then \(N=x^2+y^2+z^2=|A^\circ|^2\) and the three Weyl eigenvalues
are \((x^2-N/3)/2,(y^2-N/3)/2,(z^2-N/3)/2\).
Consequently the pointwise deficit satisfies the sharp inequality
\begin{equation}\label{eq:shape-deficit}
 D:=\tfrac16|A^\circ|^4-|W^\pm|^2_{\rm HS}
   =\tfrac12(x^2y^2+x^2z^2+y^2z^2)
   \ \ge\ \tfrac89(-\det A)_+ .
\end{equation}
For negative determinant, reverse normal if necessary, permute,
and rescale the principal curvatures by a positive constant to
write them as \((-s,1-s,1-s+v,1-s+v+u)\), with \(0<s<1\),
\(v,u\ge0\). Put \(z_0=1+2v+u\), \(b=1-2s\).
Then \(32D=2z_0^2(1+u^2)+(1-u^2)^2\ge2z_0^2+1\), whereas
\(16(-\det A)=(1-b^2)((z_0+b)^2-u^2)\). Cauchy--Schwarz gives
\[
 (1-b^2)(z_0+b)^2\le(2z_0^2+1)(1-b^2)(1/2+b^2)
                         \le\tfrac9{16}(2z_0^2+1).
\]
This proves \eqref{eq:shape-deficit}; \((-1,3,3,3)\) realizes
equality. For nonnegative determinant the formula for \(D\)
suffices. Integrating on a nonstandard candidate yields
\[
                E_+(g)\le B^2/6-64\pi^2/27 .
\]

For \(B\le(16\sqrt2/3)\pi\) the upper bound is \(64\pi^2/9\),
which contradicts the strict middle-index estimate. This proves
the third recognition criterion.

The same calculation gives a separate sign test for an actual
metric constructed using this representation. Use the sharp submanifold Sobolev inequality in
codimension two \cite[Corollary 2 and \S2]{Brendle2021Isoperimetric}
on \(F\) composed with \(\mathbb R^5\subset\mathbb R^6\).
The proof applies on the immersed domain: its contact function
takes a minimum there, and the area formula for the normal-bundle
map counts the preimages of each ambient point. The extra normal
has zero second fundamental form. Writing \(\kappa=|B^4|=\pi^2/2\),
the inequality for \(w=|f|^3\) and Cauchy--Schwarz give
\[
 4\kappa^{1/4}\|w\|_{4/3}\le\int_M\sqrt{|dw|^2+H^2w^2}\,dV_g,
 \qquad 16\sqrt\kappa\,\|f\|_4^2\le9\|df\|_2^2+\|Hf\|_2^2 .
\]
Thus the closed conformal-Laplacian form is at least
\((32\sqrt\kappa/3-B)\|f\|_4^2+\|Hf\|_2^2/12\).
At \(B=32\sqrt\kappa/3\), apply this to the positive first
eigenfunction. The mean curvature is nonzero somewhere since the
coordinate functions of a closed minimal immersion into Euclidean
space would be harmonic constants. This proves strict positivity
also at that endpoint. For a larger \(B\), one can test the smaller penalty
\(\|(|A^\circ|^2-H^2/12)_+\|_2\) against \(32\sqrt\kappa/3\)
with strict inequality. On a nonstandard candidate this would give
positive-class entry whenever also \(B<(4\sqrt{35}/3)\pi\), by
the integrated deficit estimate. All bounds refer to the same immersion.
\end{proof}

The immersion theorem supplies the first derivatives needed for
an immersion. The quantitative tests use its second derivatives
over the whole domain. Constructing an immersion satisfying either
test on an arbitrary smooth type remains an existence question.
